\documentclass[11pt,a4paper]{article}
% =====================================================================
%  Gemeinsame Vorlage der Arbeitsblätter
%  "Objektorientierte Programmierung mit Java" – Informatik 10 (NTG)
%  Einbinden in jedem Blatt direkt nach \documentclass: % =====================================================================
%  Gemeinsame Vorlage der Arbeitsblätter
%  "Objektorientierte Programmierung mit Java" – Informatik 10 (NTG)
%  Einbinden in jedem Blatt direkt nach \documentclass: % =====================================================================
%  Gemeinsame Vorlage der Arbeitsblätter
%  "Objektorientierte Programmierung mit Java" – Informatik 10 (NTG)
%  Einbinden in jedem Blatt direkt nach \documentclass: \input{ab-vorlage}
%  Übersetzen mit XeLaTeX, LuaLaTeX, pdfLaTeX oder Tectonic.
% =====================================================================
\usepackage[a4paper,left=18mm,right=18mm,top=26mm,bottom=17mm,headheight=15mm,headsep=4mm,footskip=8mm]{geometry}
\usepackage{iftex}
\ifPDFTeX
  \usepackage[T1]{fontenc}
  \usepackage[utf8]{inputenc}
  \usepackage{helvet}
  \usepackage[scaled=0.86]{beramono}
  \renewcommand{\familydefault}{\sfdefault}
\else
  \usepackage{fontspec}
  \setmainfont{texgyreheros}[Extension=.otf,UprightFont=*-regular,BoldFont=*-bold,ItalicFont=*-italic,BoldItalicFont=*-bolditalic]
  \setmonofont{DejaVuSansMono}[Extension=.ttf,UprightFont=*,BoldFont=*-Bold,ItalicFont=*-Oblique,BoldItalicFont=*-BoldOblique,Scale=0.86]
\fi
\usepackage[ngerman,shorthands=off]{babel}
\usepackage{xcolor,graphicx,tabularx,array,enumitem,fancyhdr,tikz,listings,multicol,calc,etoolbox}
\usepackage[most]{tcolorbox}
\usetikzlibrary{arrows.meta,positioning,calc}
\usepackage[hidelinks]{hyperref}
\setlength{\parindent}{0pt}
\setlength{\parskip}{3pt}
\setlist{nosep,leftmargin=*}

% ---------- Kopf- und Fußzeile (einheitliche Form) ----------
\newcommand{\lehrkraft}{Hau}
\newcommand{\abnummer}{}
\pagestyle{fancy}\fancyhf{}
\renewcommand{\headrulewidth}{0.4pt}
\fancyhead[L]{\small Klasse 10\\Datum:}
\fancyhead[C]{\small Themenbereich:\\2. Objektorientierte Programmierung}
\fancyhead[R]{\small \lehrkraft}
\fancyfoot[C]{\scriptsize edu-mrh.de\,·\,Informatik 10\,·\,Blatt \abnummer\,·\,Seite \thepage}
\newcommand{\abtitel}[2]{\renewcommand{\abnummer}{#1}\begin{center}\Large\bfseries #1\quad\underline{#2}\end{center}\vspace{-1mm}}

% ---------- Boxen ----------
\newenvironment{aufgabe}[2][]{\begin{tcolorbox}[enhanced,breakable,colback=white,colframe=black,boxrule=0.7pt,arc=3mm,left=2mm,right=2mm,top=1.5mm,bottom=1.5mm,before skip=2.5mm,after skip=2.5mm]\textbf{Aufgabe #2)}\ifstrempty{#1}{}{ \textit{#1}}\par\vspace{0.6mm}}{\end{tcolorbox}}
\newtcolorbox{merke}[1][Merke]{enhanced,breakable,colback=green!5,colframe=green!40!black,boxrule=0.8pt,arc=2mm,left=2mm,right=2mm,top=1.5mm,bottom=1.5mm,before skip=2.5mm,after skip=2.5mm,title={#1},fonttitle=\bfseries,coltitle=white}
\newtcolorbox{hinweis}{enhanced,breakable,colback=yellow!12,colframe=orange!70!black,boxrule=0.6pt,arc=2mm,left=2mm,right=2mm,top=1mm,bottom=1mm,before skip=2mm,after skip=2mm}
\newcommand{\web}[1]{{\small\color{gray}\textit{Website: #1}}}

% ---------- Lücken, Linien, Karos ----------
\newcommand{\luecke}[1][3.5cm]{\rule[-0.5ex]{#1}{0.4pt}}
\newcommand{\linien}[1]{\par\vspace{1.5mm}\foreach \i in {1,...,#1}{\noindent\rule[0pt]{\linewidth}{0.3pt}\par\vspace{5.5mm}}\vspace{-3mm}}
\newcommand{\kariert}[1]{\par\vspace{1.5mm}\noindent\begin{tikzpicture}\draw[step=5mm,gray!55,thin] (0,0) grid (\linewidth-1pt,#1*5mm);\end{tikzpicture}\par}
\newcommand{\karobox}[2]{\begin{tikzpicture}\draw[step=5mm,gray!55,thin] (0,0) grid (#1,#2);\end{tikzpicture}}
\newcommand{\klassenkarte}[3]{\begin{tabular}{|l|}\hline\textbf{#1}\\\hline #2\\\hline #3\\\hline\end{tabular}}
\newcommand{\code}[1]{\texttt{#1}}

% ---------- Java-Listings ----------
\lstset{language=Java,basicstyle=\ttfamily\small,keywordstyle=\bfseries,commentstyle=\color{gray!80!black},stringstyle=\color{black},showstringspaces=false,columns=flexible,keepspaces=true,tabsize=3,frame=single,framerule=0.4pt,rulecolor=\color{black},xleftmargin=2mm,framexleftmargin=1mm,aboveskip=2mm,belowskip=2mm,breaklines=true,
 literate={ä}{{\"a}}1 {ö}{{\"o}}1 {ü}{{\"u}}1 {Ä}{{\"A}}1 {Ö}{{\"O}}1 {Ü}{{\"U}}1 {ß}{{\ss}}1 {–}{{--}}1 {„}{{\glqq}}1 {“}{{\grqq}}1}
\lstnewenvironment{java}[1][]{\lstset{#1}}{}

%  Übersetzen mit XeLaTeX, LuaLaTeX, pdfLaTeX oder Tectonic.
% =====================================================================
\usepackage[a4paper,left=18mm,right=18mm,top=26mm,bottom=17mm,headheight=15mm,headsep=4mm,footskip=8mm]{geometry}
\usepackage{iftex}
\ifPDFTeX
  \usepackage[T1]{fontenc}
  \usepackage[utf8]{inputenc}
  \usepackage{helvet}
  \usepackage[scaled=0.86]{beramono}
  \renewcommand{\familydefault}{\sfdefault}
\else
  \usepackage{fontspec}
  \setmainfont{texgyreheros}[Extension=.otf,UprightFont=*-regular,BoldFont=*-bold,ItalicFont=*-italic,BoldItalicFont=*-bolditalic]
  \setmonofont{DejaVuSansMono}[Extension=.ttf,UprightFont=*,BoldFont=*-Bold,ItalicFont=*-Oblique,BoldItalicFont=*-BoldOblique,Scale=0.86]
\fi
\usepackage[ngerman,shorthands=off]{babel}
\usepackage{xcolor,graphicx,tabularx,array,enumitem,fancyhdr,tikz,listings,multicol,calc,etoolbox}
\usepackage[most]{tcolorbox}
\usetikzlibrary{arrows.meta,positioning,calc}
\usepackage[hidelinks]{hyperref}
\setlength{\parindent}{0pt}
\setlength{\parskip}{3pt}
\setlist{nosep,leftmargin=*}

% ---------- Kopf- und Fußzeile (einheitliche Form) ----------
\newcommand{\lehrkraft}{Hau}
\newcommand{\abnummer}{}
\pagestyle{fancy}\fancyhf{}
\renewcommand{\headrulewidth}{0.4pt}
\fancyhead[L]{\small Klasse 10\\Datum:}
\fancyhead[C]{\small Themenbereich:\\2. Objektorientierte Programmierung}
\fancyhead[R]{\small \lehrkraft}
\fancyfoot[C]{\scriptsize edu-mrh.de\,·\,Informatik 10\,·\,Blatt \abnummer\,·\,Seite \thepage}
\newcommand{\abtitel}[2]{\renewcommand{\abnummer}{#1}\begin{center}\Large\bfseries #1\quad\underline{#2}\end{center}\vspace{-1mm}}

% ---------- Boxen ----------
\newenvironment{aufgabe}[2][]{\begin{tcolorbox}[enhanced,breakable,colback=white,colframe=black,boxrule=0.7pt,arc=3mm,left=2mm,right=2mm,top=1.5mm,bottom=1.5mm,before skip=2.5mm,after skip=2.5mm]\textbf{Aufgabe #2)}\ifstrempty{#1}{}{ \textit{#1}}\par\vspace{0.6mm}}{\end{tcolorbox}}
\newtcolorbox{merke}[1][Merke]{enhanced,breakable,colback=green!5,colframe=green!40!black,boxrule=0.8pt,arc=2mm,left=2mm,right=2mm,top=1.5mm,bottom=1.5mm,before skip=2.5mm,after skip=2.5mm,title={#1},fonttitle=\bfseries,coltitle=white}
\newtcolorbox{hinweis}{enhanced,breakable,colback=yellow!12,colframe=orange!70!black,boxrule=0.6pt,arc=2mm,left=2mm,right=2mm,top=1mm,bottom=1mm,before skip=2mm,after skip=2mm}
\newcommand{\web}[1]{{\small\color{gray}\textit{Website: #1}}}

% ---------- Lücken, Linien, Karos ----------
\newcommand{\luecke}[1][3.5cm]{\rule[-0.5ex]{#1}{0.4pt}}
\newcommand{\linien}[1]{\par\vspace{1.5mm}\foreach \i in {1,...,#1}{\noindent\rule[0pt]{\linewidth}{0.3pt}\par\vspace{5.5mm}}\vspace{-3mm}}
\newcommand{\kariert}[1]{\par\vspace{1.5mm}\noindent\begin{tikzpicture}\draw[step=5mm,gray!55,thin] (0,0) grid (\linewidth-1pt,#1*5mm);\end{tikzpicture}\par}
\newcommand{\karobox}[2]{\begin{tikzpicture}\draw[step=5mm,gray!55,thin] (0,0) grid (#1,#2);\end{tikzpicture}}
\newcommand{\klassenkarte}[3]{\begin{tabular}{|l|}\hline\textbf{#1}\\\hline #2\\\hline #3\\\hline\end{tabular}}
\newcommand{\code}[1]{\texttt{#1}}

% ---------- Java-Listings ----------
\lstset{language=Java,basicstyle=\ttfamily\small,keywordstyle=\bfseries,commentstyle=\color{gray!80!black},stringstyle=\color{black},showstringspaces=false,columns=flexible,keepspaces=true,tabsize=3,frame=single,framerule=0.4pt,rulecolor=\color{black},xleftmargin=2mm,framexleftmargin=1mm,aboveskip=2mm,belowskip=2mm,breaklines=true,
 literate={ä}{{\"a}}1 {ö}{{\"o}}1 {ü}{{\"u}}1 {Ä}{{\"A}}1 {Ö}{{\"O}}1 {Ü}{{\"U}}1 {ß}{{\ss}}1 {–}{{--}}1 {„}{{\glqq}}1 {“}{{\grqq}}1}
\lstnewenvironment{java}[1][]{\lstset{#1}}{}

%  Übersetzen mit XeLaTeX, LuaLaTeX, pdfLaTeX oder Tectonic.
% =====================================================================
\usepackage[a4paper,left=18mm,right=18mm,top=26mm,bottom=17mm,headheight=15mm,headsep=4mm,footskip=8mm]{geometry}
\usepackage{iftex}
\ifPDFTeX
  \usepackage[T1]{fontenc}
  \usepackage[utf8]{inputenc}
  \usepackage{helvet}
  \usepackage[scaled=0.86]{beramono}
  \renewcommand{\familydefault}{\sfdefault}
\else
  \usepackage{fontspec}
  \setmainfont{texgyreheros}[Extension=.otf,UprightFont=*-regular,BoldFont=*-bold,ItalicFont=*-italic,BoldItalicFont=*-bolditalic]
  \setmonofont{DejaVuSansMono}[Extension=.ttf,UprightFont=*,BoldFont=*-Bold,ItalicFont=*-Oblique,BoldItalicFont=*-BoldOblique,Scale=0.86]
\fi
\usepackage[ngerman,shorthands=off]{babel}
\usepackage{xcolor,graphicx,tabularx,array,enumitem,fancyhdr,tikz,listings,multicol,calc,etoolbox}
\usepackage[most]{tcolorbox}
\usetikzlibrary{arrows.meta,positioning,calc}
\usepackage[hidelinks]{hyperref}
\setlength{\parindent}{0pt}
\setlength{\parskip}{3pt}
\setlist{nosep,leftmargin=*}

% ---------- Kopf- und Fußzeile (einheitliche Form) ----------
\newcommand{\lehrkraft}{Hau}
\newcommand{\abnummer}{}
\pagestyle{fancy}\fancyhf{}
\renewcommand{\headrulewidth}{0.4pt}
\fancyhead[L]{\small Klasse 10\\Datum:}
\fancyhead[C]{\small Themenbereich:\\2. Objektorientierte Programmierung}
\fancyhead[R]{\small \lehrkraft}
\fancyfoot[C]{\scriptsize edu-mrh.de\,·\,Informatik 10\,·\,Blatt \abnummer\,·\,Seite \thepage}
\newcommand{\abtitel}[2]{\renewcommand{\abnummer}{#1}\begin{center}\Large\bfseries #1\quad\underline{#2}\end{center}\vspace{-1mm}}

% ---------- Boxen ----------
\newenvironment{aufgabe}[2][]{\begin{tcolorbox}[enhanced,breakable,colback=white,colframe=black,boxrule=0.7pt,arc=3mm,left=2mm,right=2mm,top=1.5mm,bottom=1.5mm,before skip=2.5mm,after skip=2.5mm]\textbf{Aufgabe #2)}\ifstrempty{#1}{}{ \textit{#1}}\par\vspace{0.6mm}}{\end{tcolorbox}}
\newtcolorbox{merke}[1][Merke]{enhanced,breakable,colback=green!5,colframe=green!40!black,boxrule=0.8pt,arc=2mm,left=2mm,right=2mm,top=1.5mm,bottom=1.5mm,before skip=2.5mm,after skip=2.5mm,title={#1},fonttitle=\bfseries,coltitle=white}
\newtcolorbox{hinweis}{enhanced,breakable,colback=yellow!12,colframe=orange!70!black,boxrule=0.6pt,arc=2mm,left=2mm,right=2mm,top=1mm,bottom=1mm,before skip=2mm,after skip=2mm}
\newcommand{\web}[1]{{\small\color{gray}\textit{Website: #1}}}

% ---------- Lücken, Linien, Karos ----------
\newcommand{\luecke}[1][3.5cm]{\rule[-0.5ex]{#1}{0.4pt}}
\newcommand{\linien}[1]{\par\vspace{1.5mm}\foreach \i in {1,...,#1}{\noindent\rule[0pt]{\linewidth}{0.3pt}\par\vspace{5.5mm}}\vspace{-3mm}}
\newcommand{\kariert}[1]{\par\vspace{1.5mm}\noindent\begin{tikzpicture}\draw[step=5mm,gray!55,thin] (0,0) grid (\linewidth-1pt,#1*5mm);\end{tikzpicture}\par}
\newcommand{\karobox}[2]{\begin{tikzpicture}\draw[step=5mm,gray!55,thin] (0,0) grid (#1,#2);\end{tikzpicture}}
\newcommand{\klassenkarte}[3]{\begin{tabular}{|l|}\hline\textbf{#1}\\\hline #2\\\hline #3\\\hline\end{tabular}}
\newcommand{\code}[1]{\texttt{#1}}

% ---------- Java-Listings ----------
\lstset{language=Java,basicstyle=\ttfamily\small,keywordstyle=\bfseries,commentstyle=\color{gray!80!black},stringstyle=\color{black},showstringspaces=false,columns=flexible,keepspaces=true,tabsize=3,frame=single,framerule=0.4pt,rulecolor=\color{black},xleftmargin=2mm,framexleftmargin=1mm,aboveskip=2mm,belowskip=2mm,breaklines=true,
 literate={ä}{{\"a}}1 {ö}{{\"o}}1 {ü}{{\"u}}1 {Ä}{{\"A}}1 {Ö}{{\"O}}1 {Ü}{{\"U}}1 {ß}{{\ss}}1 {–}{{--}}1 {„}{{\glqq}}1 {“}{{\grqq}}1}
\lstnewenvironment{java}[1][]{\lstset{#1}}{}

\begin{document}
\abtitel{2.5}{Verwaltung gleichartiger Objekte -- Das Array / Feld}

\begin{aufgabe}[Projekt Zufallszahlen]{1}
\textbf{a)} Ergänze die Attribute \code{zahl1} bis \code{zahl3}.\quad \textbf{b)} Weise ihnen im Konstruktor Werte von 1 bis 6 zu.\\
\textbf{c)} Und bei 100 Zahlen? Skizziere, wie ein Speicher aussehen müsste, in dem du „die i-te Zahl“ ansprechen kannst:
\linien{2}
\end{aufgabe}

\begin{minipage}[t]{0.7\linewidth}
\begin{merke}[Felder]
Felder (engl. \luecke[2.5cm]) dienen der Speicherung einer \luecke[2.5cm] \luecke[2.5cm] von Werten \luecke[4cm], die über einen Index (Positionszahl) adressiert werden können.

Dadurch kann z.\,B. mit einer Wiederholung auf alle Elemente zugegriffen werden, ohne jedes einzeln zu behandeln.
\end{merke}
\end{minipage}\hfill
\begin{minipage}[t]{0.27\linewidth}
\vspace{3mm}
\begin{tabular}{|c|c|}\hline
Index i & \code{textFeld[i]}\\\hline
0 & \code{"Hallo "}\\\hline
1 & \code{"Welt, "}\\\hline
2 & \code{"ein "}\\\hline
3 & \code{"Satz."}\\\hline
\end{tabular}
\end{minipage}

\begin{aufgabe}{2}
Übertrage den Code in das Projekt (TODO 2) und beschreibe die Funktion der Zeilen.
\end{aufgabe}
\begin{java}
class Zufallszahlen {
   private int[] zahlenFeld;              // ________________________________________

   Zufallszahlen() {
      zahlenFeld = new int[8];            // ________________________________________
      zahlenFeld[0] = 5;                  // ________________________________________
      zahlenFeld[7] = Random.randint(1, 6); // ______________________________________
   }                                      // ________________________________________
}
\end{java}

\begin{minipage}[t]{0.7\linewidth}
\begin{merke}
\begin{itemize}
\item Felder beginnen ihren Index immer bei \luecke[1.2cm].
\item Die Größe eines Feldes liefert \code{<Feldname>.length}. Hier also: \luecke[3.5cm]. Das Array ist \luecke[1cm] Zellen groß.
\item Was passiert bei \code{zahlenFeld[8]}? \luecke[4.5cm]
\end{itemize}
\end{merke}
Fülle die Tabelle und erkläre, warum wir \code{zahlenFeld[7]} verwenden müssen, um auf die letzte Zelle zuzugreifen.
\linien{2}
Auf die letzte Zelle eines beliebig großen Feldes greift man also mit \luecke[5cm] zu.
\end{minipage}\hfill
\begin{minipage}[t]{0.27\linewidth}
\vspace{3mm}
\begin{tabular}{|c|c|}\hline
Index i & \code{zahlenFeld[i]}\\\hline
0 & 5\\\hline
1 & 0\\\hline
2 & \\\hline
\rule{0pt}{4.5mm} & \\\hline
\rule{0pt}{4.5mm} & \\\hline
\rule{0pt}{4.5mm} & \\\hline
\rule{0pt}{4.5mm} & \\\hline
\rule{0pt}{4.5mm} & \\\hline
\end{tabular}
\end{minipage}

\newpage
\begin{aufgabe}{3}
\textbf{a)} Weise allen Zellen zufällige Werte zu (Wiederholung mit fester Anzahl).\quad
\textbf{b)} 100 Zellen, jede mit der Summe zweier Würfel.\\
Was änderst du an der Schleife aus a), damit sie bei jeder Feldgröße stimmt?\\ \luecke[9cm]
\end{aufgabe}

\textbf{1:n-Beziehung als Sammlung umsetzen -- Arrays mit Referenzen}

\begin{center}
\begin{tikzpicture}[x=1cm,y=1cm]
\node[draw,fill=cyan!15,minimum height=6mm] (m) at (0,0) {Mannschaft};
\node[draw,fill=cyan!15,minimum height=6mm] (s) at (5,0) {Spieler};
\draw (m) -- (s) node[midway,above] {\small $<$ spielt in} node[pos=0.08,below] {\small 1} node[pos=0.92,below] {\small n};
\foreach \i/\n in {0/paul,1/tim,2/maja,3/tina} {
  \node[draw,fill=red!10,minimum width=13mm,minimum height=5mm] (z\i) at (6.6+\i*2.2,0.3) {\small\bfseries \i};
  \node[draw,fill=red!10,font=\scriptsize\ttfamily,inner sep=2pt] (o\i) at (6.6+\i*2.2,-1.1) {\n: Spieler};
  \draw[-{Stealth}] (z\i.south) -- (o\i.north);
}
\end{tikzpicture}
\end{center}
Auch Referenzen können in Arrays verwaltet werden. Dabei verweist jede Zelle auf ein \luecke[2.5cm] der anderen \luecke[2.5cm], hier: \luecke[2.5cm].

\begin{hinweis}
\textbf{Beachte:} Alle Elemente eines Feldes müssen denselben \luecke[3cm] besitzen, dieser kann aber beliebig gewählt werden: auch \code{boolean}, \code{String}, \code{Rectangle}, \ldots
\end{hinweis}

\begin{merke}[Syntaxübersicht]\small
\begin{tabular}{@{}l l@{}}
\code{<Sichtbarkeit> <Datentyp>[] <Feldname>;} & Deklaration: \code{[]} nach dem Datentyp macht ein Feld daraus.\\
\code{<Feldname> = new <Datentyp>[<Groesse>];} & Erzeugung mit fester Größe (nicht veränderbar).\\
\code{<Feldname>[<Index>]} & Zugriff auf die Zelle an dieser Position.\\
\end{tabular}

\textbf{Beispiele:}\\
lesen: \code{int eineZahl = zahlenFeld[0];}\hfill schreiben: \code{zahlenFeld[0] = eineZahl;}\\
Methode aufrufen: \code{kreisFeld[0].tuEtwas();}\hfill Attribut: \code{kreisFeld[0].xPosition = 100;}
\end{merke}

\begin{minipage}[t]{0.62\linewidth}
\begin{merke}[Mehrdimensionale Arrays]
Arrays können mehrere Dimensionen haben und so eine \luecke[3cm] aufspannen.\\
Deklaration und Erzeugung: \luecke[5cm]\\
Zugriff auf Spalte x, Zeile y: \luecke[3cm]
\end{merke}
\end{minipage}\hfill
\begin{minipage}[t]{0.34\linewidth}
\vspace{2mm}
\begin{tikzpicture}[x=6mm,y=6mm]
\fill[gray!60] (-0.4,-0.4) rectangle (5.4,5.4);
\foreach \x in {0,...,4} \foreach \y in {0,...,4} {
  \pgfmathsetmacro{\r}{int(random(0,100))}\pgfmathsetmacro{\g}{int(random(0,100))}\pgfmathsetmacro{\b}{int(random(0,100))}
  \fill[red!\r!green!\g!blue!\b] (\x,\y) rectangle (\x+0.95,\y+0.95);
}
\fill[white] (1,3) rectangle (1.95,3.95); \fill[white] (3,3) rectangle (3.95,3.95);
\fill[white] (1,1) rectangle (1.95,1.95); \fill[white] (2,0) rectangle (2.95,0.95); \fill[white] (3,1) rectangle (3.95,1.95);
\end{tikzpicture}
\end{minipage}

\begin{aufgabe}[Projekt Mosaik]{4}
\textbf{a)} TODO 1 und 2 in \code{Mosaik1}: zweidimensionales Array passender Größe erzeugen und in \code{neuesMosaik()} füllen.\\
\textbf{b)} Im Hauptprogramm auf \code{Mosaik2} wechseln und TODO 3 bis 5 bearbeiten (Smiley).\\
Welche Zellen \code{fuellung[x][y]} hast du für den Smiley gefärbt?\\ \luecke[9cm]
\end{aufgabe}
\end{document}
